\documentclass[a4paper,11pt]{book}
%\documentclass[a4paper,twoside,11pt,titlepage]{book}
\usepackage{listings}
\usepackage[utf8]{inputenc}
\usepackage[spanish]{babel}

% \usepackage[style=list, number=none]{glossary} %
%\usepackage{titlesec}
%\usepackage{pailatino}

\decimalpoint
\usepackage{dcolumn}
\newcolumntype{.}{D{.}{\esperiod}{-1}}
\makeatletter
\addto\shorthandsspanish{\let\esperiod\es@period@code}
\makeatother


%\usepackage[chapter]{algorithm}
\RequirePackage{verbatim}
%\RequirePackage[Glenn]{fncychap}
\usepackage{fancyhdr}
\usepackage{graphicx}
\usepackage{afterpage}

\usepackage{longtable}

\usepackage[pdfborder={000}]{hyperref} %referencia

% ********************************************************************
% Re-usable information
% ********************************************************************
\newcommand{\myTitle}{Título del proyecto\xspace}
\newcommand{\myDegree}{Grado en ...\xspace}
\newcommand{\myName}{Nombre Apllido1 Apellido2 (alumno)\xspace}
\newcommand{\myProf}{Nombre Apllido1 Apellido2 (tutor1)\xspace}
\newcommand{\myOtherProf}{Nombre Apllido1 Apellido2 (tutor2)\xspace}
%\newcommand{\mySupervisor}{Put name here\xspace}
\newcommand{\myFaculty}{Escuela Técnica Superior de Ingenierías Informática y de
Telecomunicación\xspace}
\newcommand{\myFacultyShort}{E.T.S. de Ingenierías Informática y de
Telecomunicación\xspace}
\newcommand{\myDepartment}{Departamento de ...\xspace}
\newcommand{\myUni}{\protect{Universidad de Granada}\xspace}
\newcommand{\myLocation}{Granada\xspace}
\newcommand{\myTime}{\today\xspace}
\newcommand{\myVersion}{Version 0.1\xspace}


\hypersetup{
pdfauthor = {\myName (email (en) ugr (punto) es)},
pdftitle = {\myTitle},
pdfsubject = {},
pdfkeywords = {palabra_clave1, palabra_clave2, palabra_clave3, ...},
pdfcreator = {LaTeX con el paquete ....},
pdfproducer = {pdflatex}
}

%\hyphenation{}


%\usepackage{doxygen/doxygen}
%\usepackage{pdfpages}
\usepackage{url}
\usepackage{colortbl,longtable}
\usepackage[stable]{footmisc}
%\usepackage{index}

%\makeindex
%\usepackage[style=long, cols=2,border=plain,toc=true,number=none]{glossary}
% \makeglossary

% Definición de comandos que me son tiles:
%\renewcommand{\indexname}{Índice alfabético}
%\renewcommand{\glossaryname}{Glosario}

\pagestyle{fancy}
\fancyhf{}
\fancyhead[LO]{\leftmark}
\fancyhead[RE]{\rightmark}
\fancyhead[RO,LE]{\textbf{\thepage}}
\renewcommand{\chaptermark}[1]{\markboth{\textbf{#1}}{}}
\renewcommand{\sectionmark}[1]{\markright{\textbf{\thesection. #1}}}

\setlength{\headheight}{1.5\headheight}

\newcommand{\HRule}{\rule{\linewidth}{0.5mm}}
%Definimos los tipos teorema, ejemplo y definición podremos usar estos tipos
%simplemente poniendo \begin{teorema} \end{teorema} ...
\newtheorem{teorema}{Teorema}[chapter]
\newtheorem{ejemplo}{Ejemplo}[chapter]
\newtheorem{definicion}{Definición}[chapter]

\definecolor{gray97}{gray}{.97}
\definecolor{gray75}{gray}{.75}
\definecolor{gray45}{gray}{.45}
\definecolor{gray30}{gray}{.94}

\lstset{ frame=Ltb,
     framerule=0.5pt,
     aboveskip=0.5cm,
     framextopmargin=3pt,
     framexbottommargin=3pt,
     framexleftmargin=0.1cm,
     framesep=0pt,
     rulesep=.4pt,
     backgroundcolor=\color{gray97},
     rulesepcolor=\color{black},
     %
     stringstyle=\ttfamily,
     showstringspaces = false,
     basicstyle=\scriptsize\ttfamily,
     commentstyle=\color{gray45},
     keywordstyle=\bfseries,
     %
     numbers=left,
     numbersep=6pt,
     numberstyle=\tiny,
     numberfirstline = false,
     breaklines=true,
   }
 
% minimizar fragmentado de listados
\lstnewenvironment{listing}[1][]
   {\lstset{#1}\pagebreak[0]}{\pagebreak[0]}

\lstdefinestyle{CodigoC}
   {
	basicstyle=\scriptsize,
	frame=single,
	language=C,
	numbers=left
   }
\lstdefinestyle{CodigoC++}
   {
	basicstyle=\small,
	frame=single,
	backgroundcolor=\color{gray30},
	language=C++,
	numbers=left
   }

 
\lstdefinestyle{Consola}
   {basicstyle=\scriptsize\bf\ttfamily,
    backgroundcolor=\color{gray30},
    frame=single,
    numbers=none
   }


\newcommand{\bigrule}{\titlerule[0.5mm]}


%Para conseguir que en las páginas en blanco no ponga cabecerass
\makeatletter
\def\clearpage{%
  \ifvmode
    \ifnum \@dbltopnum =\m@ne
      \ifdim \pagetotal <\topskip
        \hbox{}
      \fi
    \fi
  \fi
  \newpage
  \thispagestyle{empty}
  \write\m@ne{}
  \vbox{}
  \penalty -\@Mi
}
\makeatother

\usepackage{pdfpages}
\begin{document}
\begin{titlepage}
 
 
    \newlength{\centeroffset}
    \setlength{\centeroffset}{-0.5\oddsidemargin}
    \addtolength{\centeroffset}{0.5\evensidemargin}
    \thispagestyle{empty}
    
    \noindent\hspace*{\centeroffset}\begin{minipage}{\textwidth}
    
    \centering
    \includegraphics[width=0.9\textwidth]{imagenes/logo_ugr.jpg}\\[1.4cm]
    
    \textsc{ \Large TRABAJO FIN DE GRADO\\[0.2cm]}
    \textsc{ INGENIERÍA INFORMÁTICA}\\[1cm]
    % Upper part of the page
    % 
    % Title
    {\Huge\bfseries {\myTitle}\\
    }
    \noindent\rule[-1ex]{\textwidth}{3pt}\\[3.5ex]
    {\large\bfseries Un estudio un detalle sobre Curvas Ellípticas aplicadas a la Criptografía}
    \end{minipage}
    
    \vspace{2.5cm}
    \noindent\hspace*{\centeroffset}\begin{minipage}{\textwidth}
    \centering
    
    \textbf{Autor}\\ {Leon Elliott Fuller}\\[2.5ex]
    \textbf{Directores}\\
    {Jesús García Miranda\\[2cm]}
    \includegraphics[width=0.3\textwidth]{imagenes/etsiit_logo.png}\\[0.1cm]
    \textsc{Escuela Técnica Superior de Ingenierías Informática y de Telecomunicación}\\
    \textsc{---}\\
    Granada, Abril de 2025
    \end{minipage}
    %\addtolength{\textwidth}{\centeroffset}
    %\vspace{\stretch{2}}
    \end{titlepage}
    
    
    
\chapter*{}
%\thispagestyle{empty}
%\cleardoublepage

%\thispagestyle{empty}

\begin{titlepage}
 
 
    \setlength{\centeroffset}{-0.5\oddsidemargin}
    \addtolength{\centeroffset}{0.5\evensidemargin}
    \thispagestyle{empty}
    
    \noindent\hspace*{\centeroffset}\begin{minipage}{\textwidth}
    
    \centering
    %\includegraphics[width=0.9\textwidth]{imagenes/logo_ugr.jpg}\\[1.4cm]
    
    %\textsc{ \Large PROYECTO FIN DE CARRERA\\[0.2cm]}
    %\textsc{ INGENIERÍA EN INFORMÁTICA}\\[1cm]
    % Upper part of the page
    % 
    
     \vspace{3.3cm}
    
    %si el proyecto tiene logo poner aquí
    \includegraphics{imagenes/logo.png} 
     \vspace{0.5cm}
    
    % Title
    
    {\Huge\bfseries \myTitle\\}
    \noindent\rule[-1ex]{\textwidth}{3pt}\\[3.5ex]
    {\large\bfseries Un estudio un detalle sobre Curvas Ellípticas aplicadas a la Criptografía.\\[4cm]}
    \end{minipage}
    
    \vspace{2.5cm}
    \noindent\hspace*{\centeroffset}\begin{minipage}{\textwidth}
    \centering
    
    \textbf{Autor}\\ {Leon Elliott Fuller}\\[2.5ex]
    \textbf{Directores}\\
    Jesús García Miranda\\[2cm]
    %\includegraphics[width=0.15\textwidth]{imagenes/tstc.png}\\[0.1cm]
    %\textsc{Departamento de Teoría de la Señal, Telemática y Comunicaciones}\\
    %\textsc{---}\\
    %Granada, mes de 201
    \end{minipage}
    %\addtolength{\textwidth}{\centeroffset}
    \vspace{\stretch{2}}
    
     
    \end{titlepage}
    
    
    



\cleardoublepage
\thispagestyle{empty}

\begin{center}
{\large\bfseries Elliptic Curve Criptography: Un estudio un detalle sobre Curvas Ellípticas aplicadas a la Criptografía}\\
\end{center}
\begin{center}
Leon Elliott Fuller\\
\end{center}

%\vspace{0.7cm}
\noindent{\textbf{Palabras clave}: palabra\_clave1, palabra\_clave2, palabra\_clave3, ......}\\

\vspace{0.7cm}
\noindent{\textbf{Resumen}}\\

Poner aquí el resumen.
\cleardoublepage


\thispagestyle{empty}


\begin{center}
{\large\bfseries Project Title: Project Subtitle}\\
\end{center}
\begin{center}
First name, Family name (student)\\
\end{center}

%\vspace{0.7cm}
\noindent{\textbf{Keywords}: Keyword1, Keyword2, Keyword3, ....}\\

\vspace{0.7cm}
\noindent{\textbf{Abstract}}\\

Write here the abstract in English.

\chapter*{}
\thispagestyle{empty}

\noindent\rule[-1ex]{\textwidth}{2pt}\\[4.5ex]

Yo, \textbf{Leon Elliott Fuller}, alumno de la titulación Ingeniería Informática de la \textbf{Escuela Técnica Superior
de Ingenierías Informática y de Telecomunicación de la Universidad de Granada}, con NIE Y3299550F, autorizo la
ubicación de la siguiente copia de mi Trabajo Fin de Grado en la biblioteca del centro para que pueda ser
consultada por las personas que lo deseen. Todo el código
fuente así como este documento en formato LaTeX se puede encontrar en
el siguiente repositorios de GitHub

\vspace{6cm}

\noindent Fdo: Leon Elliott Fuller

\vspace{2cm}

\begin{flushright}
Granada a X de mayo de 2025.
\end{flushright}


\chapter*{}
\thispagestyle{empty}

\noindent\rule[-1ex]{\textwidth}{2pt}\\[4.5ex]

D. \textbf{Jesús García Miranda}, Profesor del Área de XXXX del Departamento YYYY de la Universidad de Granada.

\vspace{0.5cm}

\vspace{0.5cm}

\textbf{Informa:}

\vspace{0.5cm}

Que el presente trabajo, titulado \textit{\textbf{Título del proyecto, Subtítulo del proyecto}},
ha sido realizado bajo su supervisión por \textbf{Leon Elliott Fuller}, y autorizamos la defensa de dicho trabajo ante el tribunal
que corresponda.

\vspace{0.5cm}

Y para que conste, expiden y firman el presente informe en Granada a X de mayo de 2025.

\vspace{1cm}

\textbf{El tutor:}

\vspace{5cm}

\noindent \textbf{Jesús García Miranda}

\chapter*{Agradecimientos}
\thispagestyle{empty}

       \vspace{1cm}


Poner aquí agradecimientos...


%\frontmatter
%\tableofcontents
%\listoffigures
%\listoftables
%
%\mainmatter
%\setlength{\parskip}{5pt}

%\chapter{Introducción}

\section{Contexto histórico y motivación}
La comunicación confidencial entre sujetos ha sido una necesidad constante a lo largo de la historia. Cuando los mensajes debían transmitirse de forma no presencial o transportarse físicamente, surgía el reto de asegurar que ninguna entidad externa pudiera leerlos o alterarlos sin ser detectada.

Ya en la Antigüedad, civilizaciones como la espartana emplearon la “Escítala” en el siglo V a. C., un dispositivo de transposición que cifraba mensajes al enrollarlos en un cilindro. Más tarde, los griegos documentaron con Polibio un sistema de sustitución basado en el orden de las letras, y Roma adoptó el famoso Cifrado César, que desplazaba el alfabeto para ocultar información militar y diplomática.

Durante el Renacimiento, el cifrado avanzó hacia métodos más sofisticados, como el de Leon Battista Alberti que introdujo en 1465 la idea del cifrado polialfabético mediante discos concéntricos, lo que supuso un gran avance en esa época. Blaise de Vigenère perfeccionó esta técnica en el siglo XVI con una tabla que combinaba alfabetos de forma periódica, y en Europa se publicaron tratados (como el \emph{Cryptomenytices et Cryptographiae} de Selenus) que difundieron estos métodos entre cortes y ejércitos.

El siglo XX presenció el mayor salto tecnológico, donde se partió de los cifrados manuales de la Primera Guerra Mundial hasta la aparición de máquinas electromecánicas como la Enigma alemana en la Segunda Guerra Mundial, una máquina de rotores que automatizaba los cálculos que era necesario realizar para las operaciones de cifrado y descifrado de mensajes. El descifrado de Enigma por los criptógrafos aliados no solo cambió el curso del conflicto, sino que marcó el inicio de la era del cálculo automatizado al servicio de la criptografía. Tras la guerra, Claude Shannon sentó las bases teóricas de la información y la seguridad, y en los años 70 el NIST publicó el Estándar de Cifrado de Datos (DES), consolidando el cifrado simétrico en la informática moderna.

Sin embargo, todos los criptosistemas anteriores tenían un inconveniente, y es que todos se basaban en el hecho de que tanto la parte que descifraba como la que codificaba debían conocer el método de cifrado y la clave para descifrarlo. ¿Pero cómo se transmite la clave de forma segura? Por supuesto, con cifrado. Pero entonces, ¿cómo se puede enviar el cifrado, para que la entidad que descifra el mensaje pueda conocer dicho cifrado? Como se puede apreciar, dicho problema de la criptografía es un bucle infinito entre el cifrado y la distribución de claves que no se puede resolver con seguridad garantizada para todos los lectores deseados.

Hoy día, RSA y sistemas análogos sustentan protocolos críticos—HTTPS, correo firmado, servicios bancarios y más—pero presentan dos limitaciones crecientes:  
\begin{itemize}  
  \item \textbf{Tamaño de clave versus seguridad.} Para alcanzar un nivel de protección comparable al de cifrados simétricos de 128 bits, RSA requiere claves de varios miles de bits, lo que penaliza el rendimiento en cómputo y ancho de banda.  
  \item \textbf{Vulnerabilidad cuántica.} 
  Los algoritmos cuánticos plantean un desafío fundamental para la criptografía moderna. En particular, **Shor** demostró en 1994 que un ordenador cuántico suficientemente grande podría factorizar enteros y resolver problemas de logaritmo discreto en tiempo polinómico, lo que rompería sistemas como RSA, Diffie–Hellman e incluso ECC \cite{turn2search1}\cite{turn2search2}. Esta perspectiva ha impulsado la búsqueda de algoritmos resistentes a estos ataques cuánticos, así como la adopción de curvas elípticas de mayor seguridad y longitudes de clave incrementadas.  
\end{itemize}  

Estas limitaciones han impulsado la adopción de la \emph{Criptografía de Curvas Elípticas (ECC)}. Gracias a la dificultad del problema del logaritmo discreto en curvas elípticas, ECC ofrece un nivel de seguridad equivalente con claves drásticamente más pequeñas (por ejemplo, 256 bits en ECC frente a 3072 bits en RSA), mejorando así el rendimiento y la eficiencia. Además, la resistencia actual de ECC frente a los ataques cuánticos es mayor que la de RSA bajo el mismo modelo cuántico, lo que la posiciona como la opción de próxima generación para sistemas criptográficos modernos.

\section{Objetivos}
% Aquí, más adelante...

\section{Planificación}
% Aquí, más adelante...

\section{Estructura}
% Aquí, más adelante...

%
%\input{capitulos/02_EspecificacionRequisitos}
%
%\input{capitulos/03_Planificacion}
%
%\input{capitulos/04_Analisis}
%
%\input{capitulos/05_Diseno}
%
%\input{capitulos/06_Implementacion}
%
%\input{capitulos/07_Pruebas}
%
%\input{capitulos/08_Conclusiones}
%
%%\chapter{Conclusiones y Trabajos Futuros}
%
%
%%\nocite{*}
%\bibliography{bibliografia/bibliografia}\addcontentsline{toc}{chapter}{Bibliografía}
%\bibliographystyle{miunsrturl}
%
%\appendix
%\input{apendices/manual_usuario/manual_usuario}
%%\input{apendices/paper/paper}
%\input{glosario/entradas_glosario}
% \addcontentsline{toc}{chapter}{Glosario}
% \printglossary
\chapter*{}
\thispagestyle{empty}

\end{document}
